\chapter{What Phase 10 changed}\label{ch:phase10}
\chapterbrief
{the open problems and conditional statements of all earlier phases, reviewed in Phase~10 by three research sessions (phases 1--4, phase 5, phase 8), each followed by an independent audit and, where needed, a second round.}
{an honest edition has to say which earlier statements were refuted, which were closed, and which remain open. Phase~10 refuted five of them and closed many.}
{in case of a difference with Chapters~\ref{ch:interaction}, \ref{ch:regular} and \ref{ch:nonlin} this chapter prevails. The proofs of the two corrections that change statements are Dim.~A52 and A53; the other items are list cards (Part~IV).}

\section{Statements that turned out false}

\begin{proposition}{The interaction no-go with constant terms \led{S10A-T1, P1-Q6} \proved}\label{thm:s10at1}
Chapter~\ref{ch:interaction} left open (P1-Q6) whether $u_0\mapsto E_{\rm int}$ fails to be continuous in $H^m$ for every operator with at least two non-proportional terms. The statement as written is \textbf{false}: for reaction--diffusion $u_t=Du_{xx}+ru$ (two non-proportional terms) $E_{\rm int}$ is continuous from $H^1$, with
\[
E_{\rm int}(u)=(e^{sr}-1)\big(v(t)-v_0\big),\qquad s=t-t_0 .
\]
The correct form is: for monomial terms of distinct degrees, $E_{\rm int}$ is unbounded $H^m\to L^2$ for every $m$ and base time iff there are \emph{at least two terms of degree $\ge1$}; with at most one the interaction is bounded for every $m\ge1$. Constant terms only multiply the solution by a bounded exponential and neither create nor destroy the no-go.
\dimref{S10A-T1}
\end{proposition}

\begin{proposition}{Holes: the sum of kernels is closed but not dense \led{S10A-P1} \proved}\label{prop:s10ap1}
On the annulus $A=\{1<|z|<2\}$ the Laplacian is a product of the two coprime factors $\partial_t\pm i\partial_x$. The harmonic function $\log|z|$ is not in the closure of the sum of holomorphic and anti-holomorphic functions in any $H^m$, and for $m\ge1$
\[
\ker_{H^m(A)}\Delta=(\mathcal H_1+\mathcal H_2)\oplus\C\log|z| :
\]
the sum is closed, of codimension $b_1=1$, and not dense. Hence the density question of Chapter~\ref{ch:regular} has a negative answer on non-convex regions in general (P2-Q1 is false as a universal statement), and this is a \emph{differential} example of ``closed but not dense''.
\dimref{S10A-P1}
\end{proposition}

\begin{intuition}
A hole gives a period: the flux of a harmonic function through a loop around the hole. The flux is a continuous functional that vanishes on holomorphic and anti-holomorphic pieces (their derivative has a single-valued primitive), but not on $\log|z|$.
\end{intuition}

Three more statements were refuted: the extension of the convex density theorem to non-convex regions (S10B-C1, by a hole with an interior); P2-Q1 as a universal statement (for $\Delta$ on every Lipschitz region with $b_1\ge1$); and, in the variable-coefficient phase, ``$P11(e)$ with $B$ merely continuous'' and ``the pole depends only on $A(\xi_0)$ and the factorisation''.

\section{Problems closed}

\begin{itemize}
\item \textbf{Density of the sum of kernels} on every simply connected Lipschitz region, and for $\delta(P)=0$ on every Lipschitz region, modulo two cited facts [F1], [F2]; the source of the collar field was replaced by a proof \led{S10A-T3, S10COL}.
\item \textbf{C12-P7}: ``closed but not dense'' realised by the annulus.
\item \textbf{C3-Q1}, the decidability question, closed as a \emph{classification by data model} (composition of black-box impossibility, bound, analytic and jet models); no new mathematics.
\item \textbf{OPF5C-T5, INT-T1(B), S10B-T2}: the multi-soliton is a Hirota solution; the reflectionless Jost computation holds.
\item \textbf{Uniform shock rate with several births} (OPA-Q1): modulo the equivalence Oleinik $\Leftrightarrow$ Kru\v zkov [F1].
\item \textbf{Variable coefficients}: VAR-P11(e) by Volterra and Gronwall; a solution with $V=0$, $|Z|>0$ and $\mathrm{int}\,Z=\emptyset$ (Tychonoff and a Cantor-type set); Floquet in the quadratic class ($\mathrm{tr}M=-26/5$); a monic criterion in the Weyl algebra.
\item \textbf{Systems}: Maxwell has an identifying field; at fixed real time no finite product of Dirac flows with polynomial times; the pole criterion $E_0\cap F_0\ne0$.
\item \textbf{Dictionary}: the $L^1$ rate for general completely monotone memory kernels; the minimax rate $\Theta(\varepsilon^{1/3})$.
\item \textbf{Computer-assisted}: $N\le4$, and, after the triage, $N=5$ and $\delta_6(12)=7$ for the syzygy degrees (modulo a novelty hypothesis (Nov)); $\delta_6(11)=8$ with a written proof.
\end{itemize}

\section{Textbook examples and formalisation}

Five textbook examples were checked against the whole theory: free Dirac $1+1$, the harmonic oscillator, the P\"oschl--Teller potential, the fractional heat equation and Maxwell in vacuum; $80$ of $80$ checks pass, and three new arguments were proved. The formalisation in Lean~4 with Mathlib reached $37$ entries: $27$ are faithful to the ledger statement, $7$ are weaker and $3$ are different, none is vacuous, and all $85$ \texttt{\#print axioms} use only the standard axioms. Mathlib has no Weyl algebra and no $C_0$-semigroups, so the analytic parts are not formalised.

\begin{takeaway}
The corrections of Phase~10 all push in one direction: the negative results of Phases~1--3 are sharper than first stated (constants do not break the no-go; holes do break density), and the closures are explicit about what they depend on.
\end{takeaway}
